\chapter{\kesimpulan}
\label{bab:6}
This chapter presents the summary of our research on enhancing the GraphRAG framework with comprehensive dataset generation, specialized model fine-tuning, and multi-agent architecture. Additionally, we outline several promising future research directions to further improve the system.

\section{Conclusion}
In this research, we have developed MEGA-FrOG, a Multi-agent Enhanced Generation and Adaptation Framework of Open GraphRAG, which significantly improves upon the original FrOG framework through three key innovations. First, we implemented comprehensive dataset generation methodologies using both template-based and random-walk approaches. Second, we applied specialized model fine-tuning using Quantized Low-Rank Adaptation (QLoRA) techniques. Third, we designed a sophisticated multi-agent architecture with intelligent routing and adaptive strategy selection.

Our extensive evaluation across four knowledge graph domains (Wikidata, Curriculum KG, Legal KG, and GESIS KG) demonstrates substantial improvements over the original framework. The Qwen2.5-Coder-7B model achieved the highest average performance after fine-tuning (0.4812), with particularly strong results on the Curriculum KG (0.6111). Our findings indicate that knowledge graph characteristics, particularly URI representation schemes and knowledge graph size, significantly influence performance. Knowledge graphs with natural language-based URIs, such as the Curriculum KG, enable models to achieve substantially higher performance compared to those with numeric identifiers like Wikidata.

Addressing our research questions directly:

Regarding our first research question on comprehensive dataset generation methodologies, we successfully developed and implemented both template-based and random-walk approaches. The template-based approach focused on creating linguistically diverse questions with clear semantics, while the random-walk approach ensured comprehensive structural coverage of the knowledge graph. Our evaluation revealed that template-based generation produced higher-quality training data, with models fine-tuned on template datasets consistently outperforming those trained on random-walk datasets across all knowledge graphs. This finding suggests that linguistic diversity and clear semantics are more important than structural coverage for effective SPARQL generation. We recommend a hybrid approach that leverages template-based generation as the primary method, supplemented with random-walk generation to capture diverse structural patterns.

For our second research question on specialized model fine-tuning, we demonstrated that Quantized Low-Rank Adaptation (QLoRA) techniques significantly improve entity extraction and SPARQL generation accuracy compared to few-shot prompting approaches. All fine-tuned models substantially outperformed the FrOG baseline, with improvements ranging from +4.9\% on Wikidata to +33.3\% on GESIS KG. The Mistral-Nemo model showed the largest relative improvement from fine-tuning (+0.2392), while the Qwen2.5-Coder-7B model achieved the highest overall performance (0.4812). Our parameter-efficient approach updated only 2-3\% of model parameters while providing substantial improvements, making it practical for deployment scenarios with limited computational resources. Notably, our system achieved these improvements without using any few-shot examples, in contrast to FrOG which relies heavily on manually crafted few-shot prompts.

Addressing our third research question on enhancing the original pipeline with a LangGraph-based architecture, we successfully implemented a system with runtime-configurable components, expanded fallback mechanisms including Google Search integration, and comprehensive semantic logging. This enhanced architecture significantly improved adaptability through dynamic workflow decisions based on query complexity and user settings, strengthened fault tolerance through multi-layered fallback strategies, and dramatically increased operational transparency through RDF-based event tracking and process visualization. Compared to the original pipeline approach, our LangGraph implementation provides greater flexibility in handling diverse query types, more robust recovery options when primary methods fail, and unprecedented visibility into the reasoning process, enabling both real-time monitoring and retrospective analysis of system behavior.

For our fourth research question on performance across different knowledge graph domains, our evaluation demonstrated that the enhanced system significantly outperforms the original baseline across all knowledge graphs and query complexity levels. The most substantial improvements were observed on domain-specific knowledge graphs with well-defined schemas, particularly the Curriculum KG (0.6111 average performance) and GESIS KG (0.4798 average performance, with +33.3\% improvement over FrOG). Knowledge graph characteristics, especially URI representation schemes and knowledge graph size, emerged as fundamental determinants of achievable performance. The modest performance improvements on Wikidata (0.3101 average performance, +4.9\% improvement) highlight the ongoing challenges in handling extremely large and diverse knowledge graphs effectively. Our analysis also revealed that template-based fine-tuning consistently outperformed random-walk-based fine-tuning across all knowledge graph domains and dataset types, suggesting that structured, pattern-focused training provides advantages even when evaluating on diverse query types.

The objectives from Chapter 1 (\ref{sec:tujuan}) have been successfully completed through this research by:
\begin{enumerate}
    \item Developing comprehensive dataset generation methodologies using both template-based and random-walk approaches, creating diverse, high-quality NL2SPARQL training datasets across multiple knowledge graph domains that significantly improved model training and evaluation.
    
    \item Implementing specialized model fine-tuning using QLoRA techniques for both entity extraction and SPARQL generation tasks, demonstrating substantial performance improvements over the few-shot prompting approaches used in the original GraphRAG framework, with all models showing significant gains without requiring any few-shot examples.
    
    \item Designing and implementing a sophisticated multi-agent architecture using LangGraph that provides intelligent query routing, adaptive strategy selection based on query complexity, and robust fallback mechanisms, all with comprehensive transparency through RDF-based logging and process visualization.
    
    \item Conducting extensive evaluation across four knowledge graph domains and multiple query complexity levels, quantifying significant performance improvements (ranging from +4.9\% to +33.3\% over FrOG) and identifying key factors affecting system performance across different scenarios, including the critical impact of URI representation schemes and knowledge graph homogeneity.
\end{enumerate}

\section{Future Works}
Based on the limitations and findings of our current research, we identify several promising directions for future work:

\begin{enumerate}
    \item \textbf{Hybrid Dataset Generation Optimization}: Develop methods to automatically determine the optimal balance between template-based and random-walk generation for specific knowledge graph domains. This would involve creating metrics to assess dataset quality and coverage, then using these metrics to guide the generation process adaptively. Given our finding that template-based training consistently outperformed random-walk training, future work should focus on maximizing the effectiveness of template-based approaches while using random-walk methods to ensure comprehensive coverage.
    
    \item \textbf{Advanced Fine-tuning Techniques}: Future research could explore more sophisticated fine-tuning approaches beyond QLoRA, such as full model fine-tuning with gradient checkpointing or advanced parameter-efficient methods like LoRA variants, to potentially achieve even better performance. Additionally, investigating optimal hyperparameter configurations (such as the rank parameter, which we set to 64 for Wikidata/Curriculum KG and 32 for Legal/GESIS KG) could provide valuable insights for optimizing GraphRAG systems.
    
    \item \textbf{Automated Knowledge Graph Quality Assessment}: Develop tools to automatically identify and correct inconsistencies in knowledge graphs, addressing issues like those found in our evaluation datasets (particularly the QALD-9-Plus dataset issues). This could involve combining statistical analysis with reasoning-based approaches to detect anomalies and suggest corrections, ensuring more reliable benchmarks for system evaluation.
    
    \item \textbf{Federated Knowledge Graph Querying}: Extend the system to support querying across multiple knowledge graphs simultaneously, enabling more comprehensive answers that combine information from different sources. This would require developing methods for entity alignment, schema mapping, and distributed query execution, building on our findings about the importance of URI representation schemes.
    
    \item \textbf{Reinforcement Learning from User Feedback}: Develop methods to continuously improve the system based on user feedback, using reinforcement learning to adjust model weights and generation strategies based on which approaches yield the most satisfactory results for users. This could help address the inference efficiency vs. training cost trade-offs we identified.
    
    \item \textbf{Knowledge Graph Schema Induction}: Create methods to automatically induce schema information from knowledge graphs with limited or inconsistent schema definitions. This would enable the system to work more effectively with real-world knowledge graphs that may not have complete schema documentation, addressing the challenges we observed with knowledge graphs of varying sizes and complexities.
    
    \item \textbf{Addressing Scale Challenges}: Given the modest improvements on Wikidata compared to domain-specific knowledge graphs, develop specialized techniques for handling extremely large-scale knowledge graphs with numeric identifiers. This could include hierarchical approaches, domain-specific fine-tuning, or novel encoding methods for numeric URIs.
    
\end{enumerate}

These future directions would address both the immediate limitations identified in our current implementation and the broader challenges in the field of knowledge graph question answering. By pursuing these avenues, researchers can further enhance the accessibility, accuracy, and applicability of knowledge graph systems across diverse domains and use cases.